\section{Related Work}
\label{sec:related}

\subsection{Anomaly Detection in Network Traffic}
Liu~et~al.~\cite{liu2008isolation} introduced Isolation Forest (iForest),
which isolates anomalies by recursively partitioning the feature space and
measuring the average path length needed to isolate a sample.
iForest operates in $O(n \log n)$ time and is robust to high-dimensional data,
making it well-suited for flow-level features.

Recurrent autoencoder architectures have been applied to time-series anomaly
detection~\cite{malhotra2016lstm}.
An LSTM Autoencoder reconstructs normal sequences; anomalies produce large
reconstruction errors.
Unlike iForest, it captures temporal dependencies between consecutive flows,
which is valuable for detecting slow-and-low attacks such as beaconing.

\subsection{Kubernetes Security}
Falco~\cite{falco2023} provides kernel-level syscall monitoring via eBPF,
detecting runtime anomalies such as privilege escalation and unexpected
shell spawning.
OPA Gatekeeper~\cite{opa2023} enforces admission-time policies, preventing
privileged containers and host-path mounts from entering the cluster.
Istio~\cite{istio2023} provides automatic mTLS encryption and
AuthorizationPolicy-based access control at the service mesh layer.
Prior work~\cite{shamim2020xi} has demonstrated that these controls
individually reduce attack success rates, but their joint effect with
ML-based anomaly detection has not been systematically studied.

\subsection{Ensemble Methods for IDS}
Hybrid ensemble methods that combine tree-based and neural models have shown
improved F1 scores on the NSL-KDD and UNSW-NB15 benchmarks~\cite{kim2020lstm}.
Our work differs in that we target Kubernetes-native traffic features
extracted by Zeek, and integrate the ensemble into a live SIEM pipeline
rather than evaluating in an offline setting.
